\section{What Action Cannot Close}

A determinate debt can end. Restitution is made; a sentence is served; the sum ordered by the court is paid. The act matters precisely because justice has an objective term. It would be perverse to refuse payment until affection became pure.

The creature's relation to God contains determinate obligations, but it is not itself a case that can be closed. The Sunday Mass ends; the duty for that Sunday has been fulfilled. A fast ends; the food afterward is received with thanks. Yet the giver of being has not ceased to give, and the creature has not crossed into an estate independent of him. Religion belongs to justice, but no aggregate of religious acts can make Creator and creature even.

Here correct action can become a form of spiritual laundering. The person performs an obligation, then treats the performance as permission to preserve an untouched center for the self. Worship becomes the price paid to prevent God from making any further claim. The deed may remain objectively good; the presumption lies in using it to close a relation that cannot be closed.

This is not an argument for endless anxiety. A person need not invent duties beyond divine and ecclesial command, nor repeat a fulfilled obligation until it feels fulfilled. The opposite is true: because no checklist can purchase independence, the checklist need not bear that impossible weight. Do what is due; repent when it is not done; receive mercy; refuse both self-excusing presumption and scrupulous self-redemption.

The inadequacy lies not in action but in the end assigned to it. Justice can render a due act. It cannot make God no longer the creature's last end. Even perfect compliance, considered only as compliance, would leave a further question: is the good merely obeyed, or is it loved?
